%% LyX 2.4.4 created this file.  For more info, see https://www.lyx.org/.
%% Do not edit unless you really know what you are doing.
\documentclass[english,t]{beamer}
\usepackage{mathptmx}
\usepackage[T1]{fontenc}
\usepackage[latin9]{inputenc}
\usepackage{amssymb}
\usepackage{cancel}

\makeatletter

%%%%%%%%%%%%%%%%%%%%%%%%%%%%%% LyX specific LaTeX commands.
%% Because html converters don't know tabularnewline
\providecommand{\tabularnewline}{\\}

%%%%%%%%%%%%%%%%%%%%%%%%%%%%%% Textclass specific LaTeX commands.
% this default might be overridden by plain title style
\newcommand\makebeamertitle{\frame{\maketitle}}%
% (ERT) argument for the TOC
\AtBeginDocument{%
  \let\origtableofcontents=\tableofcontents
  \def\tableofcontents{\@ifnextchar[{\origtableofcontents}{\gobbletableofcontents}}
  \def\gobbletableofcontents#1{\origtableofcontents}
}

%%%%%%%%%%%%%%%%%%%%%%%%%%%%%% User specified LaTeX commands.
\usetheme{Warsaw}
% or ...

\setbeamercovered{transparent}
% or whatever (possibly just delete it)
\setbeamercovered{invisible} %makes overlays invisible

%gets rid of navigation symbols
\setbeamertemplate{navigation symbols}{}

%\usepackage{hyperref}
%\usepackage[usenames,dvipsnames,table]{xcolor} %adds additional color options
\definecolor{links}{HTML}{2A1B81}
%\hypersetup{colorlinks,linkcolor=blue,urlcolor=links}

\usepackage{multicol}

\usepackage{hyperref}

\usepackage{tikz}
\usetikzlibrary{calc}
\usetikzlibrary{plotmarks}
\usetikzlibrary{shapes}
\usepackage{pgfplots} %\usepackage{pgfplots, pgfplotstable}
\pgfplotsset{compat=newest}
\usepackage{filecontents}  %allows in file data
\usepgfplotslibrary{statistics} %allows histogram stuff

\usepackage{forest} %%for tree diagram

\usepackage{calc}
\usepackage[nomessages]{fp}

\usepackage[super]{nth} %easy 1st, 2nd, etc

%\usepackage{advdate}    % Advancing/saving dates
\usepackage{datetime}   % Dates formatting
%\usepackage{datenumber} % Counters for dates

\usepackage[super]{nth}

\usepackage{etex}

\usepackage{fancybox}

%%table colors
%\usepackage[table]{xcolor}
\usepackage{colortbl}

\usepackage[outline]{contour}  %allows text halos

\usepackage{epsdice} %%dice font
\usepackage[misc]{ifsym}
\newfont\dice{dice3d} %%3d dice font

\contourlength{1pt} %sets halo width
%%example: \node at (0.75,0.5) {\contour{white}{\Large Text!}};
%%example:  \node [fill=white,inner sep=1pt] at (2.25,0.5) {\Large 			Text!};
%%example:  \node [fill=white,rounded corners=2pt,inner sep=1pt] at 			(3.75,0.5) {\Large Text!};

%%%%%%%%%%%%%%%%%%%%%%%%%%%
\def\formatdateny{\csname noyear\languagename\expandafter\endcsname\formatdate}
\def\noyearenglish#1, \the\@year{#1}
\let\noyearamerican\noyearenglish
\let\noyearbritish\noyearenglish

%%allows uncover with forest diagrams
\tikzset{
    invisible/.style={opacity=0,text opacity=0},
    visible on/.style={alt=#1{}{invisible}},
    alt/.code args={<#1>#2#3}{%
      \alt<#1>{\pgfkeysalso{#2}}{\pgfkeysalso{#3}} % \pgfkeysalso doesn't change the path
    },
}
\forestset{
  visible on/.style={
    for tree={
      /tikz/visible on={#1},
      edge={/tikz/visible on={#1}}}}}

%%%redefines column alignment to match vertically
%\define@key{beamerframe}{t}[true]{% top
  %\beamer@frametopskip=.2cm plus .5\paperheight\relax%
  %\beamer@framebottomskip=0pt plus 1fill\relax%
 % \beamer@frametopskipautobreak=\beamer@frametopskip\relax%
  %\beamer@framebottomskipautobreak=\beamer@framebottomskip\relax%
  %\def\beamer@initfirstlineunskip{}%
%}

%%provides pdf with 4 slides per page; don't forget to add 'handout' to remove overlays
%\usepackage{pgfpages}
%\pgfpagesuselayout{4 on 1}[a4paper, landscape, border shrink=7mm]
%\pgfpageslogicalpageoptions{1}{border code=\pgfusepath{stroke}}
%\pgfpageslogicalpageoptions{2}{border code=\pgfusepath{stroke}}
%\pgfpageslogicalpageoptions{3}{border code=\pgfusepath{stroke}}
%\pgfpageslogicalpageoptions{4}{border code=\pgfusepath{stroke}}
%%%Don't forget to add 'handout'

\makeatother

\usepackage{babel}
\begin{document}
%%data for exams%%
\begin{filecontents}{exam_data.csv} 
dist 
89
85 
33 
51 
74 
40 
92 
64 
68 
92 
84 
96 
98 
60 
85 
71 
96 
24 
87 
81 
64 
71 
58 
40
\end{filecontents}

\newcommand{\xmax}{2000}
\newcommand{\xmin}{0}
\newcommand{\xstep}{0} %must be xmin+n
\newcommand{\ymax}{500}
\newcommand{\ymin}{0}
\newcommand{\ystep}{0} %must be ymin+n

\newcounter{countEx} \setcounter{countEx}{0}
\newcounter{countProb} \setcounter{countProb}{0}
\resetcounteronoverlays{countEx} \resetcounteronoverlays{countProb}

\newcommand{\ex}[3]{
	\addtocounter{countEx}{1}
	\only<#1-#2>{\begin{exampleblock} {Example \chNum.\secNum.\arabic{countEx}.} #3	\end{exampleblock}}
}

\newcommand{\prob}[4]{
	\addtocounter{countProb}{1}
	\only<#1-#2>{\begin{exampleblock} {Problem  \chNum.\secNum.\arabic{countProb}.\,#3} #4	\end{exampleblock}}
}

\newcommand{\yline}[4]{
	(#2 - #4)/(#1 - #3)*(x - #1) + #2
}

\beamerdefaultoverlayspecification{<+->}

%%%%Chapter and Section numbers%%%%%
\newcommand{\chNum}{4}
\newcommand{\secNum}{2}
\newcommand{\secTitle}{Measures of Central Tendency}
\date{\formatdate{31}{10}{2014}}

\title[Section \chNum.\secNum: \secTitle]{Section \chNum.\secNum: \secTitle}

\subtitle{Finite Mathematics~~MTH 107~~Fall 2014}

\author{Dr.\,James Ashe}

\titlegraphic{\includegraphics[height=2.75cm]{/home/jim/JamesAshe/MHU/images/MarsHilUniversityLogo-Virtical.png}}

\institute{Mars Hill University}

\makebeamertitle 

%%True false command
\newcommand{\T}{\textcolor{green}{\contour{black}{T}}}
\newcommand{\F}{\textcolor{red}{\contour{black}{F}}}
\newcommand{\uT}[2]{\uncover<#1-#2>{\T}}
\newcommand{\uF}[2]{\uncover<#1-#2>{\F}}

%tkiz ball item 
\newcommand*\circled[1]{\tikz[baseline=(char.base)]
	{\node[circle,ball color=green!50!black!100, shade,   color=white,inner sep=1.2pt] (char) {\tiny #1};
	}
}
\begin{frame}{Happy Halloween}

\begin{center}
\textbf{Data:}
\[
70,70,80,90,100
\]
\par\end{center}
\begin{columns}


\column{4cm}

\textbf{Mean}
\begin{itemize}
\item []<2->''average''
\item []\uncover<3->{$\frac{70+70+80+90+100}{5}$}
\end{itemize}

\column{4cm}

\textbf{Median}
\begin{itemize}
\item []<4->''middle''
\item []\uncover<5->{$70,70,$\framebox{80}$,90,100$}
\end{itemize}

\column{4cm}

\textbf{Mode}
\begin{itemize}
\item []<6->''most frequent''
\item []\uncover<7->{\framebox{70,70}$,80,90,100$}
\end{itemize}
\end{columns}
\end{frame}

\begin{frame}{The Mean}

\begin{columns}


\column{5.5cm}
\begin{block}{Mean Given $n$ data points, $x_{1},x_{2},\dots,x_{n}$, the \emph{mean
}is\emph{ $\bar{x}=\dfrac{\sum x}{n}$.}}
\end{block}
\begin{itemize}
\item []<4->$n=6$
\item []<5->$\sum x=3+15+22+9+17+13$
\item []<5-> \hspace{0cm}\hphantom{$\sum x$}$=79$
\item []
\item []<6->$\bar{x}=\dfrac{\sum x}{n}=\dfrac{79}{6}\approx13.167$
\end{itemize}

\column{5.5cm}

\ex{2}{}{Find the mean of the given data set.}\vspace{.25cm}

\begin{center}
\uncover<3->{%
\begin{tabular}{|c|c|}
\hline 
\rowcolor{yellow!50}\textbf{Year} & \textbf{TDs}\tabularnewline
\hline 
$1985$ & 3\tabularnewline
\hline 
$1986$ & 15\tabularnewline
\hline 
$1987$ & 22\tabularnewline
\hline 
$1988$ & 9\tabularnewline
\hline 
$1989$ & 17\tabularnewline
\hline 
$1990$ & 13\tabularnewline
\hline 
\end{tabular}}
\par\end{center}

\end{columns}
\end{frame}

\begin{frame}{Mean: Grouped Data}

\begin{columns}


\column{4.5cm}
\begin{itemize}
\item []<4->$n=277,000$
\item []<5->$\sum x=?$
\item []<6->Use the midpoint, $\frac{\mbox{bottom}+\mbox{top}}{2}$.

\begin{itemize}
\item <7->[ex.]$\frac{16+20}{2}=18$
\end{itemize}
\item []<9->also need to adjust for the frequency.

\begin{itemize}
\item [ex.]<10->$108,000\cdot18=1,944,000$
\end{itemize}
\item []<15->Mean: $\bar{x}=\frac{\sum(f\cdot x)}{n}$\uncover<16->{$=\frac{8,116,500}{277,00}$}
\end{itemize}

\column{6.75cm}

\ex{2}{}{Find the mean of the given data set.}\vspace{.25cm}\scriptsize

\begin{center}
\uncover<3->{\hspace{-.525cm}%
\begin{tabular}{|c|c|c|l|}
\hline 
\rowcolor{yellow!50}\tiny\uncover<7->{\textbf{Midpoint}} & \textbf{Wage} & \textbf{Freq.} & \uncover<11->{$f\cdot x$}\tabularnewline
\hline 
\uncover<7->{$18$} & $16\leq x<20$ & $108,000$ & \tiny\uncover<12->{$1,944,000$}\tabularnewline
\hline 
\uncover<8->{$22.5$} & $20\leq x<25$ & $53,000$ & \tiny\uncover<13->{$1,192,500$}\tabularnewline
\hline 
\uncover<8->{$30$} & $25\leq x<35$ & $41,000$ & \tiny\uncover<13->{$1,230,000$}\tabularnewline
\hline 
\uncover<8->{$40$} & $35\leq x<45$ & $23,000$ & \tiny\uncover<13->{$920,000$}\tabularnewline
\hline 
\uncover<8->{$50$} & $45\leq x<50$ & $29,000$ & \tiny\uncover<13->{$1,450,000$}\tabularnewline
\hline 
\uncover<8->{$60$} & $50\leq x<55$ & $23,000$ & \tiny\uncover<13->{$1,380,000$}\tabularnewline
\hline 
\multicolumn{1}{c}{} &  & $n=277,00$ & \tiny\uncover<14->{$\sum(f\cdot x)$}\tabularnewline
\cline{3-3}
\multicolumn{1}{c}{} & \multicolumn{1}{c}{} &  & \uncover<14->{$=8,116,500$}\tabularnewline
\cline{4-4}
\end{tabular}}
\par\end{center}

\end{columns}
\end{frame}

\begin{frame}{The Median}

\begin{columns}


\column{5cm}
\begin{block}<5->{Location of Median}
Given $n$ data points, the location of the median is $L=\dfrac{n+1}{2}$
after the data has been arranged in ascending order.
\end{block}

\column{6cm}

\ex{1}{}{Find the median of the following data:
\begin{enumerate}
\item [a.] $2,4,6,8,12,10,7$
\item [b.] $2,4,6,8,12,10$
\end{enumerate}
}
\begin{itemize}
\item <2->[a.]2,4,6,\framebox{7},8,10,12
\item <3->[b.]When $n$ is even, take the average of the middle:

\begin{itemize}
\item []<4->$2,4,$\framebox{6,8}$,12,10$
\item []<6->Location: $L=\dfrac{6+1}{2}=$$3.5$ data pt between $6$ and
$8$,
\item []<7->$\frac{6+8}{2}=7$
\end{itemize}
\end{itemize}
\end{columns}
\end{frame}

\begin{frame}{The Mode}

\begin{columns}


\column{5.5cm}
\begin{itemize}
\item <2->[a.] 1 4 5 5 8 10 10 10 

\begin{itemize}
\item []<3->The mode is 10
\end{itemize}
\item <4->[b.] 1 1 4 4 9 9 10 10 

\begin{itemize}
\item []<4->Each data point has the same frequency
\item []<5->-there is no mode.
\end{itemize}
\item <6->[c.] 1 3 3 6 8 9 9 

\begin{itemize}
\item []$2\times3's$ and $2\times9's$
\item []The mode is 3 and 9
\item []called \emph{bimodal.}
\end{itemize}
\end{itemize}

\column{5.5cm}

\ex{1}{}{Find the mode of the following data sets:
\begin{itemize}
\item [a.] 4 10 1 8 5 10 5 10
\item [b.] 4 9 1 10 1 10 4 9
\item [c.] 9 6 1 8 3 10 3 9
\end{itemize}
}
\end{columns}
\end{frame}

\end{document}
